\documentclass{report}

\usepackage{algpseudocode}
\usepackage{hyperref}

\title{Exercício 1 - Palíndromos de RNA}
\author{Mateus Regasi Gomes Martins}
\date{\today}


\begin{document}

\begin{titlepage}  
  
% faculdade no topo  
\begin{center}  
UNIVERSIDADE FEDERAL FLUMINENSE - UFF
\end{center}  
  
\vspace{8cm}  
  
% título central  
\begin{center}  
{  
EXERCÍCIO 1 - PALÍNDROMOS DE DNA
}  
\end{center}  
  
\vspace{1cm}  
  
% nome deslocado para direita  
\begin{flushright}  
MATEUS REGASI GOMES MARTINS
\end{flushright}  
  
\vfill  
  
% data no fundo  
\begin{center}  
NITEROI -- RJ\\  
2026  
\end{center}  
  
\end{titlepage}

\section{INTRODUÇÃO}
O seguinte texto é um relatório concernente a implementação de um algoritmo que identifica palíndromos de bases nitrogenadas em sequências de DNA realizado pela disciplina de Biologia Computacional.

\section{PALÍNDROMO DE DNA}
Seja $\Sigma = \{A,T,C,G\}$ e seja uma string $S \in \Sigma^*$ uma sequência de DNA arbitrária que pode ser decomposta em $S=a_1a_2\dots a_n$ onde $a_i$ são as bases nitrogenadas. Seja $C(S)$ a operação de complemento das bases nitrogenadas, de maneira que $C(S) = \overline{a_1}\overline{a_2}\dots\overline{a_n}$ e que $\overline A = T$, $\overline{T} = A$, $\overline{C} = G$ e $\overline{G} = C$. Seja $S^{-1}$ a operação de inversão na sequência de DNA, de maneira que $S^{-1} = a_na_{n-1}a_{n-2}\dots a_2 a_1$.

Uma sequência de bases nitrogenadas S é dita ser um palíndromo se, e somente se, $C(S^{-1}) = S$. 

\section{ALGORITMO}
Seja um arquivo $E$ composto por $n$ entradas de tamanho $m$, de maneira que $E = E_1, E_2, \dots, E_n$. Seja $E_i^{(j)}$ o caractere encontrado na posição $j$ da entrada $E_i$. Seja $L=[(i_0,j_0), (i_1,j_1), \dots, (i_p,j_p)]$ uma lista composta por trechos de DNA contendo posição inicial $i$ e posição final $j$ dos palíndromos. Seja $K$ o limite inferior do tamanho da cadeia dentro de $L$.

Assim, seque o seguinte algoritmo:

\begin{algorithmic}
\While{$\forall E_i \in E$}
    \State $E_i \gets$ leia a entrada.
    \State $L \gets$ cria a lista de trechos de DNA
    \State $l \gets 0$
    \While{$l \le m$}
        \If{$E_i^{l} = \overline(E_i^{l+1})$}
            \State $i \gets l$;
            \State $j \gets l+1$;
            
            \While{$i > 0 \land j < m-1 \land E_i^{i-1} = E_i^{j+1}$}
                \State $i \gets i-1$
                \State $j \gets j+1$
            \EndWhile

            \If{$j - i + 1 \geq K$}
                \State $L \gets L + (i, j)$
            \EndIf
        \EndIf
    \EndWhile
    
    \While{$\forall (i,j) \in L$}
        \If{$(i_{l-1} \leq i_l \land j_{l-1} \geq j_l) \lor (i_{l+1} \leq i_l \land j_{l+1} \geq j_l)$}
            \State $L \gets L - (i,j)$
        \EndIf
    \EndWhile

\EndWhile
\end{algorithmic}

\section{RESULTADOS}
\subsection{MAIOR PALÍNDROMO ENCONTRADO}
Dentro das sequências propostas foram encontradas diferentes palíndromos. A primeira sequência proposta foram duas sequências:

ACTATCCGCGTTTTTTCCAAAGTGAGCAAAAATGAAAGCTACGATCCCCCCCCCCTGAAGTTATATG
AGTGTTTTTGATAGAGCGTAAATATAAAACGTTTTATAGCCGTAATCGAAAAGCGCGATACTAAAAAAAAAACTATGAAAAAAAACTCTTG

Foi necessário definir $K = 91$ e completar o restante da primeira fita com espaços (para ficar com exatamente 91 caracteres). Obtem-se o seguinte trecho: TATAAAACGTTTTATA (i=22, j=37) de tamanho 16.

Já para a database proposta (Maribacter sp HTCC2170\footnote{Em: \url{https://www.ncbi.nlm.nih.gov/genbank/samplerecord}}) foi rodado o algoritmo em toda a base de dados. Encontrou-se na linha 25030 um palíndromo de tamanho 40 presente no trecho (i=19, j=58): TTATCGATGAACTACACAATATTGTGTAGTTCATCGATAA



\subsection{ENZIMAS DE RESTRIÇÃO}

Dentre as sequências encontradas no primeiro e curto dataset, foram encontradas\footnote{Encontradas em: \url{https://www.cellco.com.br/}} algumas enzimas de restrição que são capazes de dividir o DNA com essas bases nitrogenadas:

\begin{enumerate} 
	\item CGCG é cortado por Mbol e BssHII
	\item AGCT é cortado por Hind III 
	\item GATC é cortado por BamH I, Xba I, Nhe I e Mbol 
	\item GCGC é cortado por Mlul
\end{enumerate}

\subsection{Repositório}

Repositório presente no Github em: \url{https://github.com/mateusregasi/Exercicios-programacao/tree/main/biocomp/ex01/ex01.c}


\end{document}